\chapter{Freie Verbände und allgemeine Identitäten}
\label{b45guide:ch:free}

\section{Universelle Erzeugung}

\begin{definition}[Freier Verband]
\label{b45guide:def:free-lattice}
Ein Verband $F(X)$ mit einer Abbildung $\iota:X\to F(X)$ heißt
\emph{frei auf $X$}, wenn jede Abbildung $g:X\to L$ in einen
beliebigen Verband $L$ genau eine Fortsetzung zu einem
Verbandshomomorphismus $\widehat g:F(X)\to L$ besitzt:
\[
 \widehat g\circ\iota=g.
\]
Hier werden nur $\wedge,\vee$ als Operationen mitgeführt;
benannte Konstanten $0,1$ gehören nicht zur Definition.
\end{definition}

\begin{theorem}[Existenz, Eindeutigkeit und Test auf Identitäten]
\label{b45guide:thm:free-existence}
Zu jeder Menge $X$ existiert ein freier Verband $F(X)$.
Die Abbildung $\iota$ ist injektiv, und $\iota[X]$ erzeugt $F(X)$.
Jeder weitere freie Verband auf $X$ ist zu $F(X)$ eindeutig
durch einen die benannten Erzeuger erhaltenden Isomorphismus isomorph.

Für Verbandsterme $s,t$ in Variablen aus $X$ gilt
\[
 s=t\text{ in allen Verbänden}
 \quad\Longleftrightarrow\quad
 s(\iota[X])=t(\iota[X])\text{ in }F(X).
\]
Die rechte Seite meint die Auswertung jeder Variablen $x$ als
$\iota(x)$. Bei $X=\varnothing$ ist $F(X)$ nach der hier
verwendeten Trägerkonvention leer. \cite{nation-free}
\end{theorem}

\begin{theorem}[Größe freier Verbände mit endlich vielen Erzeugern]
\label{b45guide:thm:free-size}
Es sei $F(n)$ der freie Verband auf $n$ Erzeugern. Dann
\[
 |F(0)|=0,\qquad |F(1)|=1,\qquad |F(2)|=4.
\]
Für jedes endliche $n\geq3$ ist $F(n)$ abzählbar unendlich.
Damit können bereits drei Elemente eines Verbandes einen
unendlichen Unterverband erzeugen. \cite{nation-free}
\end{theorem}

\section{Whitmans Strukturkriterium}

\begin{theorem}[Whitman-Bedingung]
\label{b45guide:thm:whitman}
Für beliebige Elemente $u,v,w,z$ eines freien Verbandes gilt
\[
 u\wedge v\leq w\vee z
\]
genau dann, wenn mindestens eine der vier Bedingungen erfüllt ist:
\[
 \begin{array}{ll}
 u\leq w\vee z,&v\leq w\vee z,\\
 u\wedge v\leq w,&u\wedge v\leq z .
 \end{array}
\]
Für einen freien Erzeuger $x$ gelten zusätzlich
\[
\begin{aligned}
 x\leq u\vee v&\iff x\leq u\text{ oder }x\leq v,\\
 u\wedge v\leq x&\iff u\leq x\text{ oder }v\leq x .
\end{aligned}
\]
Diese Kriterien gelten in freien Verbänden und werden nicht
als Zusatzaxiome für die übrigen Verbände des Leitfadens verlangt.
\cite{nation-free,fjn}
\end{theorem}

\begin{theorem}[Entscheidbarkeit des freien Wortproblems]
\label{b45guide:thm:word-problem}
Es gibt ein terminierendes Verfahren, das für zwei beliebige
endliche Verbandsterme $s,t$ entscheidet, ob $s\leq t$ im
freien Verband gilt und ob $s=t$ im freien Verband gilt.
Damit ist entscheidbar, ob eine vorgelegte Termgleichung allein
aus den Verbandsaxiomen folgt, also in allen Verbänden gilt.
Die Aussage betrifft Terme in einem freien Verband; eine beliebige
zusätzlich vorgegebene Präsentation ist damit nicht gemeint.
\cite{nation-free,fjn}
\end{theorem}

\section{Endliche Modelle erkennen ungültige Identitäten}

\begin{definition}[Residual endlicher Verband]
Ein Verband $L$ heißt \emph{residual endlich}, wenn es für
jedes Paar $a\ne b$ in $L$ einen endlichen Verband $K$ und
einen Verbandshomomorphismus $h:L\to K$ mit $h(a)\ne h(b)$ gibt.
\end{definition}

\begin{theorem}[Endliche Trennung freier Verbände]
\label{b45guide:thm:residual}
Jeder freie Verband ist residual endlich.
Insbesondere gilt für jede Verbandstermgleichung $s=t$:
\[
\begin{aligned}
 &s=t\text{ gilt in allen Verbänden}\\
 &\qquad\Longleftrightarrow
 s=t\text{ gilt in allen endlichen Verbänden}.
\end{aligned}
\]
Eine ungültige Verbandsidentität besitzt also stets ein endliches
Gegenmodell. Eine feste gemeinsame Größenbeschränkung für alle
solchen Gegenmodelle wird hier nicht behauptet.
\cite{dean,freese-free}
\end{theorem}
